\documentclass[10pt,a4paper]{amsart}
\usepackage{amsmath,amssymb,amsthm,graphicx}
% Retained minimal preamble for the Stacks Project chapter flat.tex (More on Flatness).
% Verbatim excerpt of the chapter preamble: Xy-pic setup (native preamble lines 14-18) and the native
% theorem declarations with the mathematical macros (native preamble lines 154-217). Copied mechanically;
% nothing added, removed or redefined. The wrapper class and the amsmath/amssymb/amsthm/graphicx base
% packages are supplied by the assembler, so the class/externaldocument part of the native preamble is omitted.
% For commutative diagrams we use Xy-pic
\usepackage[all]{xy}

% We use 2cell for 2-commutative diagrams.
\xyoption{2cell}
% Theorem environments.
%
\theoremstyle{plain}
\newtheorem{theorem}[subsection]{Theorem}
\newtheorem{proposition}[subsection]{Proposition}
\newtheorem{lemma}[subsection]{Lemma}

\theoremstyle{definition}
\newtheorem{definition}[subsection]{Definition}
\newtheorem{example}[subsection]{Example}
\newtheorem{exercise}[subsection]{Exercise}
\newtheorem{situation}[subsection]{Situation}

\theoremstyle{remark}
\newtheorem{remark}[subsection]{Remark}
\newtheorem{remarks}[subsection]{Remarks}

\numberwithin{equation}{subsection}

% Macros
%
\def\lim{\mathop{\mathrm{lim}}\nolimits}
\def\colim{\mathop{\mathrm{colim}}\nolimits}
\def\Spec{\mathop{\mathrm{Spec}}}
\def\Hom{\mathop{\mathrm{Hom}}\nolimits}
\def\Ext{\mathop{\mathrm{Ext}}\nolimits}
\def\SheafHom{\mathop{\mathcal{H}\!\mathit{om}}\nolimits}
\def\SheafExt{\mathop{\mathcal{E}\!\mathit{xt}}\nolimits}
\def\Sch{\mathit{Sch}}
\def\Mor{\mathop{\mathrm{Mor}}\nolimits}
\def\Ob{\mathop{\mathrm{Ob}}\nolimits}
\def\Sh{\mathop{\mathit{Sh}}\nolimits}
\def\NL{\mathop{N\!L}\nolimits}
\def\CH{\mathop{\mathrm{CH}}\nolimits}
\def\proetale{{pro\text{-}\acute{e}tale}}
\def\etale{{\acute{e}tale}}
\def\QCoh{\mathit{QCoh}}
\def\Ker{\mathop{\mathrm{Ker}}}
\def\Im{\mathop{\mathrm{Im}}}
\def\Coker{\mathop{\mathrm{Coker}}}
\def\Coim{\mathop{\mathrm{Coim}}}

% Boxtimes
%
\DeclareMathSymbol{\boxtimes}{\mathbin}{AMSa}{"02}

%
% Macros for moduli stacks/spaces
%
\def\QCohstack{\mathcal{QC}\!\mathit{oh}}
\def\Cohstack{\mathcal{C}\!\mathit{oh}}
\def\Spacesstack{\mathcal{S}\!\mathit{paces}}
\def\Quotfunctor{\mathrm{Quot}}
\def\Hilbfunctor{\mathrm{Hilb}}
\def\Curvesstack{\mathcal{C}\!\mathit{urves}}
\def\Polarizedstack{\mathcal{P}\!\mathit{olarized}}
\def\Complexesstack{\mathcal{C}\!\mathit{omplexes}}
% \Pic is the operator that assigns to X its picard group, usage \Pic(X)
% \Picardstack_{X/B} denotes the Picard stack of X over B
% \Picardfunctor_{X/B} denotes the Picard functor of X over B
\def\Pic{\mathop{\mathrm{Pic}}\nolimits}
\def\Picardstack{\mathcal{P}\!\mathit{ic}}
\def\Picardfunctor{\mathrm{Pic}}
\def\Deformationcategory{\mathcal{D}\!\mathit{ef}}
\usepackage{cleveref}
\crefname{originaldefinition}{original definition}{original definitions}
\Crefname{originaldefinition}{Original definition}{Original definitions}
\makeatletter
\newlabel{example-two-generators}{{0EVH}{source}{}{sourcepointer-1}{}}
\newlabel{example-two-generators@cref}{{[originaldefinition][0][]0EVH}{[1][0][]source}{}{}{}}
\makeatother
\crefname{originalsection}{original section}{original sections}
\Crefname{originalsection}{Original section}{Original sections}
\makeatletter
\newlabel{coherent-section-cohomology-projective-space}{{01XS}{source}{}{sourcepointer-2}{}}
\newlabel{coherent-section-cohomology-projective-space@cref}{{[originalsection][0][]01XS}{[1][0][]source}{}{}{}}
\makeatother
\crefname{originallemma}{original lemma}{original lemmas}
\Crefname{originallemma}{Original lemma}{Original lemmas}
\makeatletter
\newlabel{more-morphisms-lemma-check-h1-fibre-zero}{{0E7F}{source}{}{sourcepointer-3}{}}
\newlabel{more-morphisms-lemma-check-h1-fibre-zero@cref}{{[originallemma][0][]0E7F}{[1][0][]source}{}{}{}}
\makeatother
\crefname{originallemma}{original lemma}{original lemmas}
\Crefname{originallemma}{Original lemma}{Original lemmas}
\makeatletter
\newlabel{more-morphisms-lemma-h1-fibre-zero}{{0E7G}{source}{}{sourcepointer-4}{}}
\newlabel{more-morphisms-lemma-h1-fibre-zero@cref}{{[originallemma][0][]0E7G}{[1][0][]source}{}{}{}}
\makeatother
\crefname{originallemma}{original lemma}{original lemmas}
\Crefname{originallemma}{Original lemma}{Original lemmas}
\makeatletter
\newlabel{more-morphisms-lemma-h1-fibre-zero-check-h0-kappa}{{0E7I}{source}{}{sourcepointer-5}{}}
\newlabel{more-morphisms-lemma-h1-fibre-zero-check-h0-kappa@cref}{{[originallemma][0][]0E7I}{[1][0][]source}{}{}{}}
\makeatother
\title{Acyclicity and square-zero kernel for the almost blowup square of an equation in two generators}
\author{The Stacks Project Authors}
\date{stacks-project:flat.tex:local-snapshot-20261011; GNU FDL 1.2 or later}
\begin{document}
\maketitle
\noindent\textit{Editorial note.} All delivered mathematical text is the author's own, copied from the pinned Stacks Project chapter file flat.tex; no statement or proof step was rewritten, reordered or removed.\par\medskip
\noindent\textit{Editorial note.} The preamble is the chapter's own retained minimal preamble (sources/stacks-flat/retained-preamble.tex), a verbatim excerpt of the native Xy-pic setup and of the native theorem declarations and mathematical macros; no macro was invented or redefined, and the wrapper class and the amsmath/amssymb/amsthm/graphicx base packages are supplied by the assembler.\par\medskip
\noindent\textit{Editorial note.} The counted statement and proof contain no cross reference, no citation and no included file, so the delivered body is self contained.\par\medskip
\noindent\textit{Editorial note.} Canonical location of the counted proposition: Stacks tag 0EVJ (sources/stacks-flat/pinned-original-source.tex line 11782).\par\medskip
\noindent\textit{Allowed standard prerequisites (editorial).}\begin{itemize}
\item The almost blowup square of an equation in two generators. Same chapter, flat.tex, Example labelled example-two-generators (source line 11488), invoked in the statement and at proof line 11812. Supplies the square Z, X, X', E and the description X = Spec(A), X' = Proj(A[T\_0,T\_1]/(f\_2T\_0 - f\_1T\_1)) used throughout.
\item Cohomology of the structure sheaf on a projective space over a ring. Cohomology of Schemes, Section labelled coherent-section-cohomology-projective-space (Stacks tag 01XS), invoked at proof lines 11804-11806. Supplies the vanishing and the global sections of the terms of the resolution in the universal case.
\item Fibrewise criteria for the vanishing of H\textasciicircum{}1 and for the surjectivity of H\textasciicircum{}0 after base change. More on Morphisms, Lemmas labelled more-morphisms-lemma-check-h1-fibre-zero (Stacks tag 0E7F), more-morphisms-lemma-h1-fibre-zero (Stacks tag 0E7G) and more-morphisms-lemma-h1-fibre-zero-check-h0-kappa (Stacks tag 0E7I), invoked at proof lines 11814-11817. Supply the base change argument that transfers the universal computation to an arbitrary A; the kernel computation is done by hand inside the counted span.
\end{itemize}\medskip
\ifdefined\hypertarget\hypertarget{sourcepointer-1}{}\fi
\noindent\textit{Source reference.} Original definition 0EVH: sources/stacks-flat/pinned-original-source.tex, source lines 11488--11488. Same-chapter pointer into the author environment example at source line 11488; invoked at proof line 11812 of the counted proof. It is retained only as a navigation pointer; every load-bearing step is inside the delivered span.\par\medskip
\ifdefined\hypertarget\hypertarget{sourcepointer-2}{}\fi
\noindent\textit{Source reference.} Original section 01XS: sources/stacks-chow/tags.txt, source lines 2531--2531. Cross-chapter standard prerequisite located in the project tag list staged in the local snapshot; invoked at proof line 11806 of the counted proof. It is retained only as a navigation pointer; every load-bearing step is inside the delivered span.\par\medskip
\ifdefined\hypertarget\hypertarget{sourcepointer-3}{}\fi
\noindent\textit{Source reference.} Original lemma 0E7F: sources/stacks-chow/tags.txt, source lines 17780--17780. Cross-chapter standard prerequisite located in the project tag list staged in the local snapshot; invoked at proof line 11815 of the counted proof. It is retained only as a navigation pointer; every load-bearing step is inside the delivered span.\par\medskip
\ifdefined\hypertarget\hypertarget{sourcepointer-4}{}\fi
\noindent\textit{Source reference.} Original lemma 0E7G: sources/stacks-chow/tags.txt, source lines 17781--17781. Cross-chapter standard prerequisite located in the project tag list staged in the local snapshot; invoked at proof line 11816 of the counted proof. It is retained only as a navigation pointer; every load-bearing step is inside the delivered span.\par\medskip
\ifdefined\hypertarget\hypertarget{sourcepointer-5}{}\fi
\noindent\textit{Source reference.} Original lemma 0E7I: sources/stacks-chow/tags.txt, source lines 17783--17783. Cross-chapter standard prerequisite located in the project tag list staged in the local snapshot; invoked at proof line 11817 of the counted proof. It is retained only as a navigation pointer; every load-bearing step is inside the delivered span.\par\medskip
\begin{lemma}
\label{lemma-funny-blow-up}
Let $Z, X, X', E$ be an almost blow up square as in
Example \ref{example-two-generators}.
Then $H^p(X', \mathcal{O}_{X'}) = 0$ for $p > 0$ and
$\Gamma(X, \mathcal{O}_X) \to \Gamma(X', \mathcal{O}_{X'})$
is a surjective map of rings whose kernel is an ideal of square zero.
\end{lemma}

\begin{proof}
First assume that $A = \mathbf{Z}[f_1, f_2]$ is the polynomial ring.
In this case our almost blow up square is the blowing up of $X = \Spec(A)$
in the closed subscheme $Z$ and in fact $X' \subset \mathbf{P}^1_X$
is an effective Cartier divisor cut out by the global section
$f_2T_0 - f_1 T_1$ of $\mathcal{O}_{\mathbf{P}^1_X}(1)$.
Thus we have a resolution
$$
0 \to
\mathcal{O}_{\mathbf{P}^1_X}(-1) \to
\mathcal{O}_{\mathbf{P}^1_X} \to
\mathcal{O}_{X'} \to
0
$$
Using the description of the cohomology given in
Cohomology of Schemes, Section
\ref{coherent-section-cohomology-projective-space}
it follows that in this case
$\Gamma(X, \mathcal{O}_X) \to \Gamma(X', \mathcal{O}_{X'})$
is an isomorphism and $H^1(X', \mathcal{O}_{X'}) = 0$.

\medskip\noindent
Next, we observe that any diagram as in Example \ref{example-two-generators}
is the base change of the diagram in the previous paragraph by the
ring map $\mathbf{Z}[f_1, f_2] \to A$. Hence by More on Morphisms, Lemmas
\ref{more-morphisms-lemma-check-h1-fibre-zero},
\ref{more-morphisms-lemma-h1-fibre-zero}, and
\ref{more-morphisms-lemma-h1-fibre-zero-check-h0-kappa}
we conclude that $H^1(X', \mathcal{O}_{X'})$ is zero in general
and the surjectivity of the map
$H^0(X, \mathcal{O}_X) \to H^0(X', \mathcal{O}_{X'})$ in general.

\medskip\noindent
Next, in the general case, let us study the kernel. If
$a \in A$ maps to zero, then looking on affine charts
we see that
$$
a = (f_1x - f_2)(a_0 + a_1x + \ldots + a_rx^r)\text{ in }A[x]
$$
for some $r \geq 0$ and $a_0, \ldots, a_r \in A$ and similarly
$$
a = (f_1 - f_2y)(b_0 + b_1y + \ldots + b_s y^s)\text{ in }A[y]
$$
for some $s \geq 0$ and $b_0, \ldots, b_s \in A$. This means we
have
$$
a = f_2 a_0,\ f_1 a_0 = f_2 a_1,\ \ldots,\ f_1 a_r = 0,
\ a = f_1 b_0,\ f_2 b_0 = f_1 b_1,\ \ldots,\ f_2 b_s = 0
$$
If $(a', r', a'_i, s', b'_j)$ is a second such system, then we have
$$
aa' = f_1f_2a_0b'_0 = f_1f_2a_1b'_1 = f_1f_2a_2b'_2 = \ldots = 0
$$
as desired.
\end{proof}

\end{document}
